% =============================================================================
%  FICHA TÉCNICA — ACTIVOS RESPIRATORIOS (AR 01–11)
%  Documento indexable para el Departamento de Indexación de Datos
%  Versión: 1.0 — Junio 2026
%  Proyecto: Aeter-AR
% =============================================================================

\documentclass[12pt,a4paper]{article}

% ── Paquetes ─────────────────────────────────────────────────────────────────
\usepackage[utf8]{inputenc}
\usepackage[T1]{fontenc}
\usepackage[spanish]{babel}
\usepackage{geometry}
\usepackage{xcolor}
\usepackage{booktabs}
\usepackage{longtable}
\usepackage{tabularx}
\usepackage{titlesec}
\usepackage{fancyhdr}
\usepackage{hyperref}
\usepackage{enumitem}

% ── Geometría ────────────────────────────────────────────────────────────────
\geometry{margin=2.5cm, top=3cm, bottom=3cm}

% ── Hipervínculos ────────────────────────────────────────────────────────────
\hypersetup{
    colorlinks=true,
    linkcolor=blue!60!black,
    urlcolor=blue!60!black,
    citecolor=blue!60!black
}

% ── Encabezado / Pie ─────────────────────────────────────────────────────────
\pagestyle{fancy}
\fancyhf{}
\fancyhead[L]{\small\textbf{Ficha Técnica — AR 01–11}}
\fancyhead[R]{\small Aeter-AR v1.0}
\fancyfoot[C]{\thepage}
\renewcommand{\headrulewidth}{0.4pt}

% ── Estilo de secciones ──────────────────────────────────────────────────────
\titleformat{\section}{\large\bfseries\color{blue!60!black}}{}{0pt}{}[\vspace{-0.3cm}\rule{\textwidth}{0.4pt}]
\titleformat{\subsection}{\bfseries\color{black!80!black}}{}{0pt}{}

% ── Colores personalizados ───────────────────────────────────────────────────
\definecolor{verde}{HTML}{2E7D32}
\definecolor{amarillo}{HTML}{F57F17}
\definecolor{rojo}{HTML}{C62828}
\definecolor{gris}{HTML}{616161}
\definecolor{azul}{HTML}{1565C0}

% ── Comandos ─────────────────────────────────────────────────────────────────
\newcommand{\veredicto}[1]{%
    \ifthenelse{\equal{#1}{fuerte}}{\textcolor{verde}{\textbf{✅ FUERTE}}}{%
    \ifthenelse{\equal{#1}{moderado}}{\textcolor{amarillo}{\textbf{⚠️ MODERADO}}}{%
    \ifthenelse{\equal{#1}{parcial}}{\textcolor{amarillo}{\textbf{⚠️ PARCIAL}}}{%
    \ifthenelse{\equal{#1}{sinrespaldo}}{\textcolor{rojo}{\textbf{❌ SIN RESPALDO}}}{%
    \ifthenelse{\equal{#1}{exagerado}}{\textcolor{rojo}{\textbf{❌ EXAGERADO}}}{#1}}}}}
}
\newcommand{\riesgo}[1]{%
    \ifthenelse{\equal{#1}{NULO}}{\textcolor{verde}{Nulo}}{%
    \ifthenelse{\equal{#1}{BAJO}}{\textcolor{amarillo}{Bajo}}{%
    \ifthenelse{\equal{#1}{MEDIO}}{\textcolor{rojo}{Medio}}{#1}}}
}
\newcommand{\riesgoficha}[1]{%
    \ifthenelse{\equal{#1}{NULO}}{\textcolor{verde}{\textbf{Nulo}}}{%
    \ifthenelse{\equal{#1}{BAJO}}{\textcolor{amarillo}{\textbf{Bajo}}}{%
    \ifthenelse{\equal{#1}{MEDIO}}{\textcolor{rojo}{\textbf{Medio}}}{#1}}}
}

\usepackage{ifthen}

\begin{document}

\thispagestyle{empty}
\begin{center}
    {\LARGE \textbf{Ficha Técnica}}\\[0.2cm]
    {\Large Activos Respiratorios — AR 01–11}\\[0.3cm]
    \rule{0.6\textwidth}{0.4pt}\\[0.3cm]
    {\small Proyecto \textbf{Aeter-AR} — Versión 1.0 — Junio 2026}\\[0.1cm]
    {\small Documento indexable para el Departamento de Indexación de Datos}\\[0.1cm]
    {\small Clasificación: \textbf{Uso interno} — Validación científica v1.0}
\end{center}

\vspace{0.5cm}

\tableofcontents
\newpage

% =============================================================================
\section{Resumen de Estructura}
% =============================================================================

Cada AR se presenta con la siguiente estructura indexable:

\begin{description}[itemsep=2pt]
    \item[ID] Identificador único del activo (AR-01 a AR-11).
    \item[Nombre] Nombre corto del activo.
    \item[Nombre Oficial] Nombre técnico completo.
    \item[Clasificación] Tipo de técnica: estimulación vagal, modulación autonómica, activación, etc.
    \item[Mecanismo] Descripción del mecanismo fisiológico subyacente.
    \item[Patrón] Esquema de tiempos de respiración (inhalación-retención-exhalación-retención).
    \item[Duración] Tiempo total recomendado de práctica.
    \item[Fases] Descripción detallada de cada fase del ciclo.
    \item[Métricas] Valores reportados de HRV, cortisol, tiempo pico, duración residual.
    \item[Riesgo] Nivel de riesgo: Nulo, Bajo, Medio.
    \item[Veredicto Científico] Evaluación del respaldo científico según literatura revisada por pares.
    \item[Correcciones] Diferencias entre el PDF original y los datos validados.
    \item[Contexto] Recomendación de uso situacional.
    \item[Fuente] Referencia bibliográfica principal.
\end{description}

\vspace{0.3cm}
\begin{tabularx}{\textwidth}{lX}
    \toprule
    \textbf{Abreviación} & \textbf{Significado} \\
    \midrule
    HRV & Heart Rate Variability (Variabilidad de la Frecuencia Cardíaca) \\
    SNS & Sistema Nervioso Simpático \\
    PNS & Sistema Nervioso Parasimpático \\
    HPA & Eje Hipotálamo-Pituitaria-Adrenal \\
    CO\textsubscript{2} & Dióxido de carbono \\
    O\textsubscript{2} & Oxígeno \\
    RCT & Ensayo Controlado Aleatorizado (Randomized Controlled Trial) \\
    VSB & Respiración Lenta Voluntaria (Voluntary Slow Breathing) \\
    \bottomrule
\end{tabularx}

\newpage

% =============================================================================
%  AR-01
% =============================================================================
\section{AR-01: Respiración Diafragmática}
\label{ar:01}

\subsection*{Identificación}
\begin{tabularx}{\textwidth}{lX}
    \textbf{ID} & AR-01 \\
    \textbf{Nombre} & Respiración Diafragmática \\
    \textbf{Nombre Oficial} & Respiración Abdominal \\
    \textbf{Clasificación} & Estimulación vagal forzada por presión diafragmática rítmica \\
\end{tabularx}

\subsection*{Patrón y Duración}
\begin{tabularx}{\textwidth}{lX}
    \textbf{Patrón} & 4-0-4-6 (inhalación 4s, pausa plena 0s, exhalación 4s, pausa vacía 6s) \\
    \textbf{Tiempo por ciclo} & 14 segundos \\
    \textbf{Duración total} & 3--5 minutos \\
    \textbf{Ciclos por minuto} & $\sim$4.3 rpm \\
\end{tabularx}

\subsection*{Mecanismo}
La contracción rítmica del diafragma estimula mecánicamente el nervio vago,
activando el sistema parasimpático. La pausa vacía de 6 segundos prolonga la
señal vagal, reduciendo la frecuencia cardíaca y la presión arterial. Meta-análisis
de 223 estudios (Laborde et al., 2022) confirma aumento de HRV vagal-mediada.

\subsection*{Fases del Ciclo}
\begin{tabularx}{\textwidth}{lccX}
    \toprule
    \textbf{Fase} & \textbf{Duración} & \textbf{Vía} & \textbf{Notas} \\
    \midrule
    Inhalación & 4s & Nasal & Solo expansión abdominal. El pecho no se mueve. \\
    Pausa Plena & 0s & — & Sin retención. Transición directa. \\
    Exhalación & 4s & Bucal & Abdomen se hunde suavemente. \\
    Pausa Vacía & 6s & — & Relajación total. Interruptor vagal. \\
    \bottomrule
\end{tabularx}

\subsection*{Métricas Reportadas}
\begin{tabularx}{\textwidth}{lX}
    \toprule
    \textbf{Métrica} & \textbf{Valor} \\
    \midrule
    Reducción de cortisol & Significativa (reducción documentada) \\
    HRV & Aumento moderado a alto \\
    Tiempo pico & 4 minutos \\
    Duración residual & 60--90 minutos \\
    \bottomrule
\end{tabularx}

\subsection*{Riesgo}
\riesgoficha{NULO} — Segura para todas las personas. Si hay mareo, reducir tiempos.

\subsection*{Veredicto Científico}
\veredicto{fuerte}

\subsection*{Correcciones Respecto al PDF Original}
\begin{itemize}
    \item ``$-27\%$ cortisol'' reemplazado por ``reducción significativa'' documentada.
    \item ``$+48\%$ HRV'' reemplazado por ``aumento moderado a alto''.
    \item Tiempos de pico (4 min) y residual (90 min) no verificados específicamente;
          se mantienen como rangos referenciales.
\end{itemize}

\subsection*{Contexto de Uso}
Uso general: estrés, inicio de cualquier práctica, recuperación.

\subsection*{Fuente}
Laborde S. et al. (2022). Effects of voluntary slow breathing on HRV: A systematic
review and meta-analysis. \textit{Neuroscience \& Biobehavioral Reviews}, 138, 104711.

\bigskip
\rule{\textwidth}{0.2pt}

% =============================================================================
%  AR-02
% =============================================================================
\section{AR-02: Physiological Sigh}
\label{ar:02}

\subsection*{Identificación}
\begin{tabularx}{\textwidth}{lX}
    \textbf{ID} & AR-02 \\
    \textbf{Nombre} & Physiological Sigh \\
    \textbf{Nombre Oficial} & Doble Inhalación — Reset de Pánico \\
    \textbf{Clasificación} & Reinflación alveolar + activación vagal por exhalación prolongada \\
\end{tabularx}

\subsection*{Patrón y Duración}
\begin{tabularx}{\textwidth}{lX}
    \textbf{Patrón} & 2+1 / 8--12 (inhalación 70\%, relleno, exhalación 8--12s) \\
    \textbf{Tiempo por ciclo} & $\sim$10 segundos \\
    \textbf{Duración total} & 90 segundos -- 5 minutos \\
    \textbf{Ciclos} & 3--5 repeticiones \\
\end{tabularx}

\subsection*{Mecanismo}
La doble inhalación reinfla alvéolos colapsados, maximizando el intercambio
gaseoso. La exhalación larga activa el nervio vago, reduciendo la frecuencia
cardíaca. Estudio randomizado Stanford 2023 (Balban et al.) demostró que el
\textit{cyclic sighing} superó a \textit{box breathing} y mindfulness en mejora
del ánimo y reducción de ansiedad (p $<$ 0.05).

\subsection*{Fases del Ciclo}
\begin{tabularx}{\textwidth}{lccX}
    \toprule
    \textbf{Fase} & \textbf{Duración} & \textbf{Vía} & \textbf{Notas} \\
    \midrule
    Inhalación Máxima & 3s & Nasal & Expansión pulmonar al $\sim$70\%. \\
    Inhalación Relleno & 1s & Nasal & Sorbo extra sin soltar aire previo. \\
    Exhalación Lenta & 10s (8--12) & Bucal & Continua y completa. Activa el vago. \\
    Repetición & 3--5 ciclos & — & Para efecto completo. \\
    \bottomrule
\end{tabularx}

\subsection*{Métricas Reportadas}
\begin{tabularx}{\textwidth}{lX}
    \toprule
    \textbf{Métrica} & \textbf{Valor} \\
    \midrule
    Reducción de ansiedad & Significativa vs. mindfulness (p $<$ 0.05) \\
    HRV & Aumento moderado \\
    Tiempo pico & 90 segundos \\
    Duración residual & 45 minutos \\
    \bottomrule
\end{tabularx}

\subsection*{Riesgo}
\riesgoficha{NULO} — Segura. Si la exhalación forzada genera mareo, reducir a 6--8s.

\subsection*{Veredicto Científico}
\veredicto{fuerte}

\subsection*{Correcciones Respecto al PDF Original}
\begin{itemize}
    \item Sin correcciones mayores. Se fortaleció con cita directa al estudio Stanford.
    \item Se agregó precisión sobre el tamaño del efecto y el contexto del estudio.
\end{itemize}

\subsection*{Contexto de Uso}
Ataques de pánico, ansiedad aguda, transiciones entre tareas de alta presión.

\subsection*{Fuente}
Balban M.Y. et al. (2023). Brief structured respiration practices enhance mood
and reduce physiological arousal. \textit{Cell Reports Medicine}, 4, 100895.
Stanford RCT NCT05304000.

\bigskip
\rule{\textwidth}{0.2pt}

% =============================================================================
%  AR-03
% =============================================================================
\section{AR-03: Box Breathing}
\label{ar:03}

\subsection*{Identificación}
\begin{tabularx}{\textwidth}{lX}
    \textbf{ID} & AR-03 \\
    \textbf{Nombre} & Box Breathing \\
    \textbf{Nombre Oficial} & Respiración en Caja — Estado Cero \\
    \textbf{Clasificación} & Nivelación autonómica 1:1 entre SNS y PNS \\
\end{tabularx}

\subsection*{Patrón y Duración}
\begin{tabularx}{\textwidth}{lX}
    \textbf{Patrón} & 4-4-4-4 (inhalación 4s, retención lleno 4s, exhalación 4s, retención vacío 4s) \\
    \textbf{Tiempo por ciclo} & 16 segundos \\
    \textbf{Duración total} & 4--8 minutos \\
\end{tabularx}

\subsection*{Mecanismo}
Las cuatro fases iguales equilibran el sistema nervioso autónomo. La retención con
pulmones llenos estimula barorreceptores; la retención vacía activa el vago.
Usado por fuerzas especiales (Navy SEALs) para mantener claridad bajo presión.
La retención de CO\textsubscript{2} modula quimiorreceptores, reduciendo la ansiedad.

\subsection*{Fases del Ciclo}
\begin{tabularx}{\textwidth}{lccX}
    \toprule
    \textbf{Fase} & \textbf{Duración} & \textbf{Vía} & \textbf{Notas} \\
    \midrule
    Inhalación & 4s & Nasal & Ritmo metronómico. \\
    Retención Lleno & 4s & — & Quietud absoluta. Sin tensión en glotis. \\
    Exhalación & 4s & Bucal & Flujo constante y controlado. \\
    Retención Vacío & 4s & — & Punto de equilibrio SNS/PNS. \\
    \bottomrule
\end{tabularx}

\subsection*{Métricas Reportadas}
\begin{tabularx}{\textwidth}{lX}
    \toprule
    \textbf{Métrica} & \textbf{Valor} \\
    \midrule
    Reducción de cortisol & Significativa (con práctica regular) \\
    HRV & $+55\%$ (aumento reportado, referencial) \\
    Tiempo pico & 5 minutos \\
    Duración residual & 3 horas \\
    \bottomrule
\end{tabularx}

\subsection*{Riesgo}
\riesgoficha{NULO} — Segura. No forzar las retenciones.

\subsection*{Veredicto Científico}
\veredicto{fuerte}

\subsection*{Correcciones Respecto al PDF Original}
\begin{itemize}
    \item Claim de ``elimina temblor motor fino'' ajustado a efecto reportado, no probado.
    \item $+55\%$ HRV mantenido como número referencial, no universal.
\end{itemize}

\subsection*{Contexto de Uso}
Estrés agudo, recuperación post-esfuerzo, preparación para tareas de alta concentración.

\subsection*{Fuente}
Estudios múltiples. Ver revisión Laborde et al. (2022).

\bigskip
\rule{\textwidth}{0.2pt}

% =============================================================================
%  AR-04
% =============================================================================
\section{AR-04: Coherencia Cardíaca}
\label{ar:04}

\subsection*{Identificación}
\begin{tabularx}{\textwidth}{lX}
    \textbf{ID} & AR-04 \\
    \textbf{Nombre} & Coherencia Cardíaca \\
    \textbf{Nombre Oficial} & Respiración de Resonancia \\
    \textbf{Clasificación} & Sincronización barorreceptora a 0.1 Hz \\
\end{tabularx}

\subsection*{Patrón y Duración}
\begin{tabularx}{\textwidth}{lX}
    \textbf{Patrón} & 5-0-5-0 (inhalación 5s, exhalación 5s, sin pausas) \\
    \textbf{Tiempo por ciclo} & 10 segundos \\
    \textbf{Frecuencia} & $\sim$6 respiraciones por minuto (0.10 Hz) \\
    \textbf{Duración total} & 5--15 minutos \\
\end{tabularx}

\subsection*{Mecanismo}
A $\sim$6 respiraciones/minuto, el sistema cardiovascular entra en resonancia:
frecuencia cardíaca, presión arterial y respiración se sincronizan. Esto maximiza
la amplitud de la HRV y optimiza el flujo sanguíneo a la corteza prefrontal.
Estudio con 1.8 millones de sesiones (Nature, 2025) confirma que la frecuencia
de coherencia más común es 0.10 Hz.

\subsection*{Fases del Ciclo}
\begin{tabularx}{\textwidth}{lccX}
    \toprule
    \textbf{Fase} & \textbf{Duración} & \textbf{Vía} & \textbf{Notas} \\
    \midrule
    Inhalación & 5s & Nasal & Flujo constante y sedoso. \\
    Exhalación & 5s & Nasal/Bucal & Sin pausas entre fases. \\
    \bottomrule
\end{tabularx}

\subsection*{Métricas Reportadas}
\begin{tabularx}{\textwidth}{lX}
    \toprule
    \textbf{Métrica} & \textbf{Valor} \\
    \midrule
    HRV & Aumento significativo (máximo en frecuencia de resonancia) \\
    Tiempo pico & 6 minutos \\
    Duración residual & 4--6 horas \\
    Coherencia & Óptima en 0.04--0.10 Hz \\
    \bottomrule
\end{tabularx}

\subsection*{Riesgo}
\riesgoficha{NULO} — Segura. Mantener postura erguida.

\subsection*{Veredicto Científico}
\veredicto{fuerte}

\subsection*{Correcciones Respecto al PDF Original}
\begin{itemize}
    \item ``$+310\%$ HRV'' eliminado. Reemplazado por ``aumento significativo de HRV''.
          El número no es replicable ni universal.
    \item Duración residual (4--6h) no establecida científicamente pero reportada
          por practicantes. Se mantiene como referencia.
\end{itemize}

\subsection*{Contexto de Uso}
Antes de reuniones importantes, estudio, trabajo creativo, o práctica diaria de base.

\subsection*{Fuente}
Nature Scientific Reports (2025). HRV biofeedback global study — 1.8 millones de
sesiones. Sévoz-Couche \& Laborde (2022). Neurosci Biobehav Rev, 135, 104576.

\bigskip
\rule{\textwidth}{0.2pt}

% =============================================================================
%  AR-05
% =============================================================================
\section{AR-05: 4-7-8 Modificado}
\label{ar:05}

\subsection*{Identificación}
\begin{tabularx}{\textwidth}{lX}
    \textbf{ID} & AR-05 \\
    \textbf{Nombre} & 4-7-8 Modificado \\
    \textbf{Nombre Oficial} & Inducción al Sueño \\
    \textbf{Clasificación} & Acumulación controlada de CO\textsubscript{2} para sedación neuronal \\
\end{tabularx}

\subsection*{Patrón y Duración}
\begin{tabularx}{\textwidth}{lX}
    \textbf{Patrón} & 4-7-8-0 (inhalación 4s, retención 7s, exhalación 8s, sin retención vacía) \\
    \textbf{Tiempo por ciclo} & 19 segundos \\
    \textbf{Duración total} & 8 ciclos ($\sim$2.5 minutos) \\
    \textbf{Máximo} & 8 ciclos por práctica \\
\end{tabularx}

\subsection*{Mecanismo}
La retención de 7 segundos acumula CO\textsubscript{2} en sangre, activando
quimiorreceptores que inducen una sensación de calma química. La exhalación de
8 segundos activa el vago. Desarrollado por el Dr.~Andrew Weil (Harvard).
Estudios muestran que la respiración lenta con retención puede reducir la
latencia de sueño, aunque el efecto es modesto ($\sim$55 segundos más rápido,
Duke 2017).

\subsection*{Fases del Ciclo}
\begin{tabularx}{\textwidth}{lccX}
    \toprule
    \textbf{Fase} & \textbf{Duración} & \textbf{Vía} & \textbf{Notas} \\
    \midrule
    Inhalación & 4s & Nasal & Silenciosa y suave. \\
    Retención Lleno & 7s & — & Acumulación química de CO\textsubscript{2}. \\
    Exhalación & 8s & Bucal & ``Whoosh'' suave. Relajación. \\
    Repetición & $\times$8 & — & Máximo 8 ciclos. \\
    \bottomrule
\end{tabularx}

\subsection*{Métricas Reportadas}
\begin{tabularx}{\textwidth}{lX}
    \toprule
    \textbf{Métrica} & \textbf{Valor} \\
    \midrule
    Reducción de cortisol & $-18\%$ (reportado) \\
    HRV & $+22\%$ (reportado) \\
    Tiempo pico & 12 minutos \\
    Duración residual & 8 horas \\
    \bottomrule
\end{tabularx}

\subsection*{Riesgo}
\riesgoficha{BAJO} — No forzar la retención de 7s si genera ansiedad.

\subsection*{Veredicto Científico}
\veredicto{moderado}

\subsection*{Correcciones Respecto al PDF Original}
\begin{itemize}
    \item ``Sueño en $<$12 minutos'' \textbf{eliminado}. Estudio Duke 2017 mostró
          solo 55 segundos de mejora en latencia de sueño.
    \item Reformulado como ``ayuda a conciliar el sueño'' sin prometer tiempos.
\end{itemize}

\subsection*{Contexto de Uso}
Dificultad para dormir, insomnio leve, relajación profunda antes de dormir.

\subsection*{Fuente}
Jerath R. et al. (2019). Self-Regulation of Breathing as Adjunctive Treatment of
Insomnia. \textit{Front Psychiatry}. Duke University (2017).

\bigskip
\rule{\textwidth}{0.2pt}

% =============================================================================
%  AR-06
% =============================================================================
\section{AR-06: Nadi Shodhana}
\label{ar:06}

\subsection*{Identificación}
\begin{tabularx}{\textwidth}{lX}
    \textbf{ID} & AR-06 \\
    \textbf{Nombre} & Nadi Shodhana \\
    \textbf{Nombre Oficial} & Respiración Alterna — Balance Hemisférico \\
    \textbf{Clasificación} & Modulación del ciclo nasal y balance autonómico \\
\end{tabularx}

\subsection*{Patrón y Duración}
\begin{tabularx}{\textwidth}{lX}
    \textbf{Patrón} & 4-4-4-4 alterno (inhalación, retención, exhalación, retención, con alternancia de fosa nasal) \\
    \textbf{Tiempo por ciclo} & 16 segundos \\
    \textbf{Duración total} & 6--10 minutos \\
\end{tabularx}

\subsection*{Mecanismo}
La respiración por una fosa nasal a la vez influye en el ciclo nasal ultradiano
(alternancia natural de dominancia nasal cada 2--3h). Estudios pequeños muestran
activación parasimpática y mejora en tareas de memoria espacial. El ``balance
hemisférico'' es un concepto yóguico sin respaldo neurocientífico directo, pero
la técnica tiene efectos calmantes documentados.

\subsection*{Fases del Ciclo}
\begin{tabularx}{\textwidth}{lccX}
    \toprule
    \textbf{Fase} & \textbf{Duración} & \textbf{Vía} & \textbf{Notas} \\
    \midrule
    Inhalar Izquierda & 4s & Nasal & Cerrar fosa derecha con el pulgar. \\
    Retención Lleno & 4s & — & Sincronización. \\
    Exhalar Derecha & 4s & Nasal & Cerrar fosa izquierda con anular. \\
    Retención Vacío & 4s & — & Invertir dirección y repetir. \\
    \bottomrule
\end{tabularx}

\subsection*{Métricas Reportadas}
\begin{tabularx}{\textwidth}{lX}
    \toprule
    \textbf{Métrica} & \textbf{Valor} \\
    \midrule
    HRV & $+28\%$ (reportado) \\
    Reducción de cortisol & $-24\%$ (reportado) \\
    Tiempo pico & 8 minutos \\
    Duración residual & 2 horas \\
    \bottomrule
\end{tabularx}

\subsection*{Riesgo}
\riesgoficha{NULO} — Segura.

\subsection*{Veredicto Científico}
\veredicto{parcial}

\subsection*{Correcciones Respecto al PDF Original}
\begin{itemize}
    \item ``$+22\%$ decisión'' \textbf{eliminado}. No hay estudio específico con ese número.
    \item ``Balance hemisférico'' mantenido como concepto yóguico, no como
          mecanismo neurocientífico probado.
\end{itemize}

\subsection*{Contexto de Uso}
Antes de meditar, ansiedad generalizada, transiciones mentales.

\subsection*{Fuente}
Nepal Medical College Journal — estudio de memoria espacial con Nadi Shodhana.

\bigskip
\rule{\textwidth}{0.2pt}

% =============================================================================
%  AR-07
% =============================================================================
\section{AR-07: Super-ventilación + Retención}
\label{ar:07}

\subsection*{Identificación}
\begin{tabularx}{\textwidth}{lX}
    \textbf{ID} & AR-07 \\
    \textbf{Nombre} & Super-ventilación + Retención \\
    \textbf{Nombre Oficial} & Potencia Yang — Alcalosis Controlada \\
    \textbf{Clasificación} & Alcalosis respiratoria + hipoxia controlada \\
\end{tabularx}

\subsection*{Patrón y Duración}
\begin{tabularx}{\textwidth}{lX}
    \textbf{Patrón} & 40 respiraciones rápidas + retención en vacío \\
    \textbf{Tiempo por ciclo} & Variable \\
    \textbf{Duración total} & 15 minutos \\
\end{tabularx}

\subsection*{Mecanismo}
La hiperventilación sostenida reduce el CO\textsubscript{2} sanguíneo (hipocapnia)
elevando el pH (alcalosis), lo que amortigua la señal de dolor y frío. La
retención en vacío subsiguiente genera hipoxia controlada. Los resultados
científicos son mixtos: un estudio en \textit{Scientific Reports} (2023) no
encontró diferencias significativas en HRV o bienestar (p $>$ 0.05).
Sin embargo, una revisión \textit{PLOS One} (2024) sugiere efectos
anti-inflamatorios (aumento de IL-10).

\subsection*{Fases del Ciclo}
\begin{tabularx}{\textwidth}{lccX}
    \toprule
    \textbf{Fase} & \textbf{Duración} & \textbf{Vía} & \textbf{Notas} \\
    \midrule
    Respiraciones rápidas & $\times$40 & Nasal & Inhalación 80\%, exhalación suelta. \\
    Vaciado total & — & — & Expulsar todo el aire. \\
    Retención Vacío & hasta 120s & — & NO forzar. Síntomas de alerta. \\
    Inhalación Rescate & 15s & Nasal & Restablecer equilibrio. \\
    \bottomrule
\end{tabularx}

\subsection*{Métricas Reportadas}
\begin{tabularx}{\textwidth}{lX}
    \toprule
    \textbf{Métrica} & \textbf{Valor} \\
    \midrule
    Respuesta adrenal & $+15\%$ (reportado) \\
    HRV & $-10\%$ (durante la práctica) \\
    Tiempo pico & 15 minutos \\
    Duración residual & 1 hora \\
    \bottomrule
\end{tabularx}

\subsection*{Riesgo}
\riesgoficha{MEDIO} — \textbf{Contraindicaciones:} epilepsia, cardiopatías,
embarazo, presión alta no controlada. Puede causar síncope o tetania.
Solo personas sanas y en posición segura.

\subsection*{Veredicto Científico}
\veredicto{parcial}

\subsection*{Correcciones Respecto al PDF Original}
\begin{itemize}
    \item Se agregaron advertencias de seguridad explícitas.
    \item Se documentaron los resultados mixtos de la literatura (Scientific
          Reports 2023 vs. PLOS One 2024).
    \item No recomendado para principiantes sin supervisión.
\end{itemize}

\subsection*{Contexto de Uso}
Personas entrenadas que buscan resiliencia al frío/dolor. No apta para principiantes.

\subsection*{Fuente}
Ketelhut S. et al. (2023). WHM effectiveness on cardiac autonomic function.
\textit{Scientific Reports}, 13, 17517. Almahayni O. (2024). \textit{PLOS One}, 19(3).

\bigskip
\rule{\textwidth}{0.2pt}

% =============================================================================
%  AR-08
% =============================================================================
\section{AR-08: Bio-Eficiencia}
\label{ar:08}

\subsection*{Identificación}
\begin{tabularx}{\textwidth}{lX}
    \textbf{ID} & AR-08 \\
    \textbf{Nombre} & Bio-Eficiencia \\
    \textbf{Nombre Oficial} & Optimización de Oxígeno — Efecto Bohr \\
    \textbf{Clasificación} & Elevación del umbral de CO\textsubscript{2} para activar el Efecto Bohr \\
\end{tabularx}

\subsection*{Patrón y Duración}
\begin{tabularx}{\textwidth}{lX}
    \textbf{Patrón} & 2-0-2 + Pausa Máxima (inhalación 2s, exhalación 2s, pausa máxima) \\
    \textbf{Duración total} & 10 minutos \\
\end{tabularx}

\subsection*{Mecanismo}
La respiración superficial reduce ligeramente el CO\textsubscript{2}, y la pausa
máxima hasta ``hambre de aire'' lo eleva nuevamente, desplazando la curva de
disociación de la hemoglobina (Efecto Bohr). Esto facilita la liberación de
O\textsubscript{2} a los tejidos. El mecanismo es fisiológicamente real
(Christian Bohr, 1904).

\subsection*{Fases del Ciclo}
\begin{tabularx}{\textwidth}{lccX}
    \toprule
    \textbf{Fase} & \textbf{Duración} & \textbf{Vía} & \textbf{Notas} \\
    \midrule
    Inhalación & 2s & Nasal & Superficial, bajo volumen. \\
    Exhalación & 2s & Nasal & Pasiva y relajada. \\
    Pausa Máxima & Variable & — & Hasta ``hambre de aire'' moderada. \\
    Recuperación & 3 ciclos & Nasal & Calmar inmediatamente. \\
    \bottomrule
\end{tabularx}

\subsection*{Métricas Reportadas}
\begin{tabularx}{\textwidth}{lX}
    \toprule
    \textbf{Métrica} & \textbf{Valor} \\
    \midrule
    Reducción de cortisol & $-12\%$ (reportado) \\
    HRV & $+35\%$ (reportado) \\
    Tiempo pico & 10 minutos \\
    Duración residual & 3 horas \\
    \bottomrule
\end{tabularx}

\subsection*{Riesgo}
\riesgoficha{BAJO} — La pausa máxima debe sentirse como ``hambre de aire'', no como asfixia.

\subsection*{Veredicto Científico}
\veredicto{moderado}

\subsection*{Correcciones Respecto al PDF Original}
\begin{itemize}
    \item ``$+25\%$ O\textsubscript{2} efficiency'' \textbf{eliminado}.
          No hay estudios que respalden ese número.
    \item Reformulado como ``optimiza la liberación de O\textsubscript{2}
          mediante el Efecto Bohr''.
\end{itemize}

\subsection*{Contexto de Uso}
Mejora de eficiencia respiratoria general, preparación para ejercicio aeróbico.

\subsection*{Fuente}
Christian Bohr (1904). Efecto Bohr. Mecanismo fisiológico establecido.

\bigskip
\rule{\textwidth}{0.2pt}

% =============================================================================
%  AR-09
% =============================================================================
\section{AR-09: Modo Turbo}
\label{ar:09}

\subsection*{Identificación}
\begin{tabularx}{\textwidth}{lX}
    \textbf{ID} & AR-09 \\
    \textbf{Nombre} & Modo Turbo \\
    \textbf{Nombre Oficial} & Activación Simpática Explosiva \\
    \textbf{Clasificación} & Activación del eje HPA para pico adrenérgico \\
\end{tabularx}

\subsection*{Patrón y Duración}
\begin{tabularx}{\textwidth}{lX}
    \textbf{Patrón} & 2 ciclos/segundo (exhalación explosiva + inhalación pasiva) \\
    \textbf{Tiempo por ciclo} & 0.5 segundos \\
    \textbf{Duración total} & 1--3 minutos \\
\end{tabularx}

\subsection*{Mecanismo}
La respiración rápida activa el sistema nervioso simpático y el eje HPA,
liberando adrenalina. La exhalación nasal explosiva crea presión intra-abdominal
que estabiliza el core. Similar a Bhastrika (pranayama del fuelle).

\subsection*{Fases del Ciclo}
\begin{tabularx}{\textwidth}{lccX}
    \toprule
    \textbf{Fase} & \textbf{Duración} & \textbf{Vía} & \textbf{Notas} \\
    \midrule
    Exhalación Explosiva & 0.5s & Nasal & Abdomen como pistón. \\
    Inhalación Pasiva & 0.5s & Nasal & Por descompresión. \\
    \bottomrule
\end{tabularx}

\subsection*{Métricas Reportadas}
\begin{tabularx}{\textwidth}{lX}
    \toprule
    \textbf{Métrica} & \textbf{Valor} \\
    \midrule
    Respuesta adrenal & $+45\%$ (reportado, estimado) \\
    HRV & $-25\%$ (durante la práctica) \\
    Tiempo pico & 60 segundos \\
    Duración residual & 20 minutos \\
    \bottomrule
\end{tabularx}

\subsection*{Riesgo}
\riesgoficha{BAJO} — No hacer si hay presión alta, ansiedad severa o problemas cardíacos.

\subsection*{Veredicto Científico}
\veredicto{moderado}

\subsection*{Correcciones Respecto al PDF Original}
\begin{itemize}
    \item Sin correcciones importantes. Mecanismo directo y simple.
    \item $+45\%$ adrenalina es un estimado razonable pero no verificado en
          estudio específico.
\end{itemize}

\subsection*{Contexto de Uso}
Antes de esfuerzo físico explosivo, para salir de letargo, calentamiento.

\subsection*{Fuente}
Mecanismo fisiológico estándar. Bhastrika pranayama como referencia.

\bigskip
\rule{\textwidth}{0.2pt}

% =============================================================================
%  AR-10
% =============================================================================
\section{AR-10: Cadencia Síncrona}
\label{ar:10}

\subsection*{Identificación}
\begin{tabularx}{\textwidth}{lX}
    \textbf{ID} & AR-10 \\
    \textbf{Nombre} & Cadencia Síncrona \\
    \textbf{Nombre Oficial} & Gestión de Carga — Sincronización Paso-Respiración \\
    \textbf{Clasificación} & Estabilización mecánica y acoplamiento locomotor-respiratorio \\
\end{tabularx}

\subsection*{Patrón y Duración}
\begin{tabularx}{\textwidth}{lX}
    \textbf{Patrón} & 2 pasos / 2--3 pasos (inhalar en 2 pasos, exhalar en 2--3) \\
    \textbf{Duración total} & Duración del esfuerzo \\
\end{tabularx}

\subsection*{Mecanismo}
La sincronización entre respiración y zancada optimiza la biomecánica respiratoria.
El patrón 2:2 es el más común en corredores de resistencia. El acoplamiento
locomotor-respiratorio está documentado en fisiología del ejercicio.

\subsection*{Fases del Ciclo}
\begin{tabularx}{\textwidth}{lccX}
    \toprule
    \textbf{Fase} & \textbf{Duración} & \textbf{Vía} & \textbf{Notas} \\
    \midrule
    Inhalación & 2 pasos & Nasal & Ajustable por intensidad. \\
    Exhalación & 2--3 pasos & Bucal & Flujo constante. \\
    \bottomrule
\end{tabularx}

\subsection*{Métricas Reportadas}
\begin{tabularx}{\textwidth}{lX}
    \toprule
    \textbf{Métrica} & \textbf{Valor} \\
    \midrule
    Reducción de cortisol & $-5\%$ (reportado) \\
    HRV & $+15\%$ (reportado) \\
    Tiempo pico & Inmediato \\
    Duración residual & Durante el esfuerzo \\
    \bottomrule
\end{tabularx}

\subsection*{Riesgo}
\riesgoficha{NULO} — Segura. Ajustar ratio según intensidad.

\subsection*{Veredicto Científico}
\veredicto{moderado}

\subsection*{Correcciones Respecto al PDF Original}
\begin{itemize}
    \item Claim sobre lactato ajustado a ``teórico, no probado''.
    \item Mantenido como ``práctica común en deportes de resistencia''.
\end{itemize}

\subsection*{Contexto de Uso}
Running, ciclismo, remo, o cualquier actividad de resistencia con patrón repetitivo.

\subsection*{Fuente}
Fisiología del ejercicio estándar. Acoplamiento locomotor-respiratorio documentado.

\bigskip
\rule{\textwidth}{0.2pt}

% =============================================================================
%  AR-11
% =============================================================================
\section{AR-11: Pausa de Precisión}
\label{ar:11}

\subsection*{Identificación}
\begin{tabularx}{\textwidth}{lX}
    \textbf{ID} & AR-11 \\
    \textbf{Nombre} & Pausa de Precisión \\
    \textbf{Nombre Oficial} & PRN Window — Estabilidad en Acción \\
    \textbf{Clasificación} & Bradicardia inducida para estabilidad motora \\
\end{tabularx}

\subsection*{Patrón y Duración}
\begin{tabularx}{\textwidth}{lX}
    \textbf{Patrón} & PRN (inhalación 2s, exhalación parcial 2s, pausa 3s) \\
    \textbf{Tiempo por ciclo} & 7 segundos \\
    \textbf{Uso} & Bajo demanda, no como práctica continua \\
\end{tabularx}

\subsection*{Mecanismo}
La pausa respiratoria reduce momentáneamente la frecuencia cardíaca (bradicardia)
y disminuye el temblor fisiológico, mejorando la estabilidad en tareas de
precisión. Usado en tiro, cirugía, o cualquier actividad que requiera pulso
firme. La bradicardia inducida por retención es un reflejo fisiológico conocido.

\subsection*{Fases del Ciclo}
\begin{tabularx}{\textwidth}{lccX}
    \toprule
    \textbf{Fase} & \textbf{Duración} & \textbf{Vía} & \textbf{Notas} \\
    \midrule
    Inhalación Normal & 2s & Nasal & Sin llenar al máximo. \\
    Exhalación Parcial & 2s & Bucal & 60--70\% del volumen. \\
    PAUSA PRN & 3s (máx 5s) & — & Ventana de acción. \\
    \bottomrule
\end{tabularx}

\subsection*{Métricas Reportadas}
\begin{tabularx}{\textwidth}{lX}
    \toprule
    \textbf{Métrica} & \textbf{Valor} \\
    \midrule
    HRV & $+5\%$ (reportado) \\
    Reducción de cortisol & $-10\%$ (reportado) \\
    Tiempo pico & Inmediato \\
    Duración residual & 30 segundos \\
    \bottomrule
\end{tabularx}

\subsection*{Riesgo}
\riesgoficha{NULO} — No mantener más de 5s para evitar hipoxia.

\subsection*{Veredicto Científico}
\veredicto{moderado}

\subsection*{Correcciones Respecto al PDF Original}
\begin{itemize}
    \item Sin correcciones mayores. Técnica contextual con mecanismo plausible.
    \item Aclarado que no hay estudios específicos que respalden la mejora de
          precisión en tiro, aunque es ampliamente usado en contextos tácticos.
\end{itemize}

\subsection*{Contexto de Uso}
Tiro deportivo, tareas de precisión, cirugía, fotografía con teleobjetivo.

\subsection*{Fuente}
Reflejo de buceo (\textit{diving reflex}) y bradicardia por retención respiratoria.
Fisiología estándar.

\bigskip
\rule{\textwidth}{0.2pt}

% =============================================================================
\section{Índice de Métricas Comparativas}
% =============================================================================

\begin{longtable}{lccccc}
    \toprule
    \textbf{AR} & \textbf{Cortisol} & \textbf{HRV} & \textbf{Pico} & \textbf{Residual} & \textbf{Riesgo} \\
    \midrule
    \endhead
    \bottomrule
    \endfoot
    AR-01 & Significativa $\downarrow$ & Moderado $\uparrow$ & 4 min & 60--90 min & Nulo \\
    AR-02 & — & Moderado $\uparrow$ & 90 s & 45 min & Nulo \\
    AR-03 & Significativa $\downarrow$ & $+$55\% $\uparrow$ & 5 min & 3 h & Nulo \\
    AR-04 & — & Significativo $\uparrow\uparrow$ & 6 min & 4--6 h & Nulo \\
    AR-05 & $-$18\% $\downarrow$ & $+$22\% $\uparrow$ & 12 min & 8 h & Bajo \\
    AR-06 & $-$24\% $\downarrow$ & $+$28\% $\uparrow$ & 8 min & 2 h & Nulo \\
    AR-07 & — & $-$10\% $\downarrow$* & 15 min & 1 h & Medio \\
    AR-08 & $-$12\% $\downarrow$ & $+$35\% $\uparrow$ & 10 min & 3 h & Bajo \\
    AR-09 & — & $-$25\% $\downarrow$* & 60 s & 20 min & Bajo \\
    AR-10 & $-$5\% $\downarrow$ & $+$15\% $\uparrow$ & Inmediato & Durante ej. & Nulo \\
    AR-11 & $-$10\% $\downarrow$ & $+$5\% $\uparrow$ & Inmediato & 30 s & Nulo \\
\end{longtable}
\vspace{0.2cm}
{\small * Durante la práctica. El resto son valores post-práctica.}

\end{document}
